% this file is called up by thesis.tex
% content in this file will be fed into the main document

%: ----------------------- introduction file header -----------------------
\chapter{Background}
\label{ch:background}

% the code below specifies where the figures are stored
\ifpdf
    \graphicspath{{2_background/figures/PNG/}{2_background/figures/PDF/}{2_background/figures/}}
\else
    \graphicspath{{2_background/figures/EPS/}{2_background/figures/}}
\fi


\section{Wind Turbine Condition Monitoring}
\label{sec:wt-condition-monitoring}

Wind Turbine Condition Monitoring (WTCM) is the systematic process of continuously or periodically assessing the operational health of wind turbine components, such as gearboxes, bearings, blades, and generators, through the measurement and analysis of operational parameters including vibration, temperature, acoustic emissions, and power output~\citep{li_comprehensive_2023,charabi_wind_2020,henrique_global_nodate}.
The principal objective of WTCM is the early detection of anomalies or degradation, thereby enabling timely maintenance interventions that mitigate the risk of catastrophic failure, minimise unplanned downtime, and sustain optimal energy production throughout the turbine’s service life.

Modern WTCM typically integrates dedicated Condition Monitoring Systems (CMS) with Supervisory Control and Data Acquisition (SCADA) platforms.
CMS record high-frequency sensor data related to mechanical and structural health, whereas SCADA provides low-frequency operational and environmental data, along with historical event logs including failures, repairs, and downtimes~\citep{dao_wind_2019}.
The combined datasets underpin reliability analyses, facilitate the optimisation of operational strategies, and form the foundation for data-driven maintenance planning.

Maintenance activities within the wind energy sector are generally classified into three principal categories~\citep{yang_operations_2021,henrique_global_nodate}:

Reactive (Corrective) Maintenance (CM): Interventions performed only after a failure has occurred or performance has dropped below acceptable thresholds.
While such actions restore operational capability, they are typically associated with elevated costs due to unplanned downtime, emergency repairs, and potential collateral damage to adjacent components.

Condition-Based Maintenance (CBM): A proactive strategy in which maintenance is triggered when monitored parameters, such as vibration or temperature, deviate from predefined operational thresholds.
CBM enables faults to be addressed in their early stages, reducing the likelihood of unplanned outages and extending component lifespan.

Predictive Maintenance (PdM): An advanced form of proactive maintenance that leverages statistical modeling, historical trends, and increasingly, machine learning algorithms to forecast the remaining useful life of components.
Maintenance is scheduled in advance of predicted failures, often incorporating CBM data but focusing on forecasting rather than detection.

The fundamental distinction between these strategies lies in their temporal relationship to component degradation: CBM is driven by current health indicators, PdM anticipates future failure events, and reactive maintenance responds after failures have occurred.
In practice, hybrid maintenance policies are often adopted, combining elements of each strategy to optimise cost, reliability, and risk.
Such hybrid approaches are particularly relevant in offshore wind energy applications, where logistical constraints, weather-related access limitations, and high operational costs necessitate highly efficient maintenance planning~\citep{yang_operations_2021}.


\section{Normal Behaviour Monitoring}
\label{sec:normal-behaviour-monitoring}
Normal Behaviour Monitoring (NBM) refers to the modeling of a system's nominal operational patterns in order to detect deviations that may indicate emerging faults or performance degradation.
In the context of wind turbine condition monitoring, NBM leverages historical and real-time operational data to establish a representation of “healthy” turbine behaviour.
Any statistically significant departure from this model is treated as a potential anomaly warranting further investigation or maintenance intervention.

NBM is particularly relevant for wind turbines due to their complex, non-linear dynamics and the wide variability of environmental conditions under which they operate.
Unlike threshold-based monitoring—where fixed limits on operational parameters trigger alerts—NBM captures multivariate and temporal correlations between signals, providing a more nuanced picture of system health.
However, the accuracy of NBM hinges on the availability of sufficient high-quality “normal” data and the ability of the modeling technique to cope with non-stationary operating regimes (e.g., seasonal wind patterns, load variations).

Modern Supervisory Control and Data Acquisition (SCADA) systems now generate high-frequency time-series telemetry data from various turbine components, enabling what is commonly referred to as NBM—the task of modeling normal operational patterns to detect abnormal deviations.

The methods employed in NBM broadly fall into three categories: statistical approaches, traditional (shallow) machine learning, and deep learning-based methods.

Statistical methods, such as Ordinary Least Squares (OLS), ARIMA, and Principal Component Analysis (PCA), are still used in certain settings~\citep{chesterman_overview_2023}.
These models are computationally efficient, interpretable, and data-efficient, making them suitable for constrained environments or first-pass screening.
However, their simplicity limits their ability to model complex, non-linear system behaviour.

Traditional machine learning models, including Random Forests, Gradient Boosting Machines, and Support Vector Machines, have grown more popular due to their improved performance on non-linear tasks and moderate data requirements.
These methods still require feature engineering and are sensitive to class imbalance, but they offer a strong middle ground in the performance–complexity trade-off~\citep{ben_ali_online_2018, bilendo_normal_2022, liu_wind_2022}.

Deep learning methods—especially neural network-based architectures—have rapidly become the dominant approach in recent research.
These models offer superior capability to capture non-linear dependencies in high-dimensional data but come at the cost of increased computational demand and reduced interpretability.
Within this category, Autoencoders (AEs) have emerged as a popular choice for anomaly detection, particularly in unsupervised or semi-supervised settings.
Like PCA, Autoencoders learn latent representations of normal behaviour via dimensionality reduction, but with significantly greater expressive power due to their ability to model non-linearities.

Accordingly, a growing body of research has focused on applying machine learning (ML) and deep learning (DL) techniques to enable condition monitoring and early fault detection in wind turbines.
Among these, Autoencoders have become one of the most widely adopted architectures for unsupervised and semi-supervised anomaly detection~\citep{chesterman_overview_2023}.

Recent trends point towards deep reconstructive normal behaviour models—architectures that not only learn a latent representation of healthy behaviour but also attempt to reconstruct the original signal.
By analysing reconstruction errors, these models can highlight subtle deviations that may be indicative of incipient faults.
Section~\ref{sec:deep-reconstructive} explores these models in detail, including their theoretical underpinnings, design considerations, and application to wind turbine condition monitoring.

------------------------------------

Modern Supervisory Control and Data Acquisition (SCADA) systems now generate high-frequency time-series telemetry data from various turbine components, enabling what is commonly referred to as Normal Behavior Monitoring (NBM) — the task of modeling normal operational patterns to detect abnormal deviations.

The methods employed in NBM broadly fall into three categories: statistical approaches, traditional (shallow) machine learning, and deep learning-based methods.

Statistical methods, such as Ordinary Least Squares (OLS), ARIMA, and Principal Component Analysis (PCA), are still used in certain settings~\citep{chesterman_overview_2023}.
These models are computationally efficient, interpretable, and data-efficient, making them suitable for constrained environments or first-pass screening.
However, their simplicity makes them less capable of modeling complex, non-linear system behavior.

Traditional machine learning models, including Random Forests, Gradient Boosting Machines, and Support Vector Machines, have grown more popular due to their better performance on non-linear tasks and moderate data requirements.
These methods still require some feature engineering and are sensitive to class imbalance, but they offer a strong middle ground in performance versus complexity trade-offs.
The methods are presented by~\cite{ben_ali_online_2018, bilendo_normal_2022, liu_wind_2022}.

Deep learning methods—especially neural network-based architectures—have rapidly become the dominant approach in recent research.
These models offer superior ability to capture non-linear dependencies in high-dimensional data but come at the cost of increased computational demand and reduced interpretability.
Within this category, Autoencoders (AEs) have emerged as a popular choice for anomaly detection, particularly in unsupervised or semi-supervised settings.
Just like PCA, Autoencoders learn latent representations of normal behavior via dimension reduction, but with significantly more expressive power due to their ability to model non-linearities.

Accordingly, a growing body of research has focused on applying machine learning (ML) and deep learning (DL) techniques to enable condition monitoring and early fault detection in wind turbines.
Among these, Autoencoders (AEs) have emerged as one of the most widely adopted architectures, particularly in unsupervised or semi-supervised anomaly detection settings~\citep{chesterman_overview_2023}.
